Considereu un experiment d'efecte fotoelèctric en què el càtode és una làmina de cesi que té una freqüència llindar de $4,59\times 10^{14}\ \mathrm{Hz}$.

\begin{apartat}{1,25}
Calculeu el treball d'extracció del càtode. Si il·luminem el càtode amb els diferents punters làser de la taula que hi ha a continuació, justifiqueu amb quins punters làser es produirà l'efecte fotoelèctric. Completeu la taula i representeu gràficament en la quadrícula adjunta l'energia cinètica màxima dels electrons (en eV) en funció de la freqüència dels fotons incidents (en Hz) per a un interval de freqüències entre $3,00\times 10^{14}\ \mathrm{Hz}$ i $7,00\times 10^{14}\ \mathrm{Hz}$.

\begin{center}
\resizebox{\linewidth}{!}{%
\begin{tabular}{|c|c|c|c|c|c|}
\hline
\textit{Tipus de punter làser} & \textit{Longitud d'ona (nm)} & \textit{Freqüència ($\times 10^{14}$ Hz)} & \textit{Energia fotó ($\times 10^{-19}$ J)} & \textit{$E_c$ electró ($\times 10^{-19}$ J)} & \textit{$E_c$ electró (eV)} \\
\hline
Làser blau & 460 & & & & \\
\hline
Làser verd & 532 & & & & \\
\hline
Làser infraroig & 1\,080 & & & & \\
\hline
\end{tabular}}
\end{center}

\figura[0.6]{fig1.png}
\end{apartat}

\begin{apartat}{1,25}
Il·luminem el càtode amb un làser de freqüència $8,00\times 10^{14}\ \mathrm{Hz}$. Calculeu la velocitat i la longitud d'ona de De Broglie dels electrons arrencats del càtode amb la radiació d'aquest làser.
\end{apartat}

\dades{%
  $1\ \mathrm{eV} = 1,602\times 10^{-19}\ \mathrm{J}$.\\
  $m_\mathrm{e} = 9,11\times 10^{-31}\ \mathrm{kg}$.\\
  $c = 3,00\times 10^{8}\ \mathrm{m\,s^{-1}}$.\\
  $h = 6,63\times 10^{-34}\ \mathrm{J\,s}$.}
